% --- Archon physics formalization source begin ---
% archon:physics
% archon:covers PhyXMiniProblems/problem_phyx_mini_0517.lean
% archon:source-report reports/phyx_mini/problem_phyx_mini_0517.source.json

\chapter{Physics problem phyx\_mini\_0517}
\label{ch:PhyXMiniProblems_problem_phyx_mini_0517}

\paragraph{Problem source.}
\#\# Physical scenario

A hypothetical atom as shown in figure has energy levels at \textbackslash{}( 0.00 \textbackslash{}text\{ eV\} \textbackslash{}) (the ground level), \textbackslash{}( 1.00 \textbackslash{}text\{ eV\} \textbackslash{}), and \textbackslash{}( 3.00 \textbackslash{}text\{ eV\} \textbackslash{}).

\#\# Question

What is the highest frequencies of the spectral lines this atom can emit when excited?

\#\# Answer choices

- A: A: \textbackslash{}( 9.16 \textbackslash{}times 10\textasciicircum{}\{14\} \textbackslash{})

- B: B: \textbackslash{}( 2.42 \textbackslash{}times 10\textasciicircum{}\{14\} \textbackslash{})

- C: C: \textbackslash{}( 5.50 \textbackslash{}times 10\textasciicircum{}\{14\} \textbackslash{})

- D: D: \textbackslash{}( 7.25 \textbackslash{}times 10\textasciicircum{}\{14\} \textbackslash{})

\#\# Figure caption (auxiliary; use the image as primary evidence)

The image is a diagram illustrating energy transitions, specifically emission and absorption in an atomic or molecular system. 

**Components and Details:**

- **Energy Levels:**
  - Three horizontal lines represent different energy levels at 3.00 eV, 1.00 eV, and a ground level.

- **Transitions:**
  - **Emission Transitions:**
    - Illustrated with pink arrows pointing downward from higher to lower energy levels. 
    - A transition from 3.00 eV to 1.00 eV.
    - A transition from 3.00 eV directly to the ground level.
    - A transition from 1.00 eV to the ground level.
    - Energies of 2.00 eV, 3.00 eV, and 1.00 eV are marked next to these arrows.

  - **Absorption Transitions:**
    - Illustrated with blue arrows pointing upward from lower to higher energy levels.
    - A transition from the ground level to 1.00 eV.
    - A transition from the ground level directly to the 3.00 eV level.
    - Energies of 1.00 eV and 3.00 eV are marked next to these arrows.

- **Text Labels:**
  - On the left, energy levels are labeled with eV values.
  - Below the emission transitions section, it is labeled "Emission transitions."
  - Below the absorption transitions section, it is labeled "Absorption transitions."

This diagram effectively displays how energy levels and transitions work in terms of emission and absorption processes using energy values.

\#\# Dataset metadata

Quantum Phenomena; Physical Model Grounding Reasoning, Predictive Reasoning

\paragraph{Recorded answer/context.}
D: D: \textbackslash{}( 7.25 \textbackslash{}times 10\textasciicircum{}\{14\} \textbackslash{})

\paragraph{Figure/image path.}
/root/proposal\_for\_physic/science-mango/phyx\_mini\_run/phyx\_data/test\_image/517.png

\paragraph{Formalization target.}
create a compiling Lean file with sorry bodies at `PhyXMiniProblems/problem\_phyx\_mini\_0517.lean`. The Lean declarations must preserve the physical quantities, dimensions or dimensional roles, figure labels, governing-law hypotheses, and final relation expressed by this problem.
Use Mathlib/Physlib names found through LeanExplore where available. If a physics API is missing, introduce faithful local abstractions rather than scalar placeholder aliases.

\begin{theorem}[Physics formalization target]
\label{thm:physics:phyx_mini_0517:target}
\lean{PhyXMiniProblems.ProblemPhyXMini0517.highestEmittedSpectralLine_is_choiceD}
\uses{def:physics:phyx-mini-0517:phyxminiproblems-problemphyxmini0517-threelevelatomsetup, def:physics:phyx-mini-0517:phyxminiproblems-problemphyxmini0517-matchesexcitedatomscenario, def:physics:phyx-mini-0517:phyxminiproblems-problemphyxmini0517-matchessuppliedenergyleveldiagram, def:physics:phyx-mini-0517:phyxminiproblems-problemphyxmini0517-hasphysicalspectralparameters, def:physics:phyx-mini-0517:phyxminiproblems-problemphyxmini0517-usesstandardplanckconstant, def:physics:phyx-mini-0517:phyxminiproblems-problemphyxmini0517-satisfiesplanckeinsteinemissionlaw, def:physics:phyx-mini-0517:phyxminiproblems-problemphyxmini0517-isuniquehighestemissiontransition, def:physics:phyx-mini-0517:phyxminiproblems-problemphyxmini0517-agreeswithdisplayedfrequency, def:physics:phyx-mini-0517:phyxminiproblems-problemphyxmini0517-isuniqueclosestanswerchoice}
The assigned autoformalize agent should translate this physics problem into Lean declarations in the covered file, with theorem and lemma proof bodies written as `by sorry`.
\end{theorem}
\begin{proof}
This is an autoformalization task, not a proof task. Produce statements that can later be proved by the physics prover without weakening the physical model.
\end{proof}

% --- Archon declaration topology begin ---
\section{Declaration topology}
\begin{definition}[action Dimension]
  \label{def:physics:phyx-mini-0517:phyxminiproblems-problemphyxmini0517-actiondimension}
  \lean{PhyXMiniProblems.ProblemPhyXMini0517.actionDimension}
  The physical dimension of action, equivalently energy times time
\end{definition}
\begin{proof}
  This is the declared model definition of action Dimension.
\end{proof}
\begin{definition}[Action Quantity]
  \label{def:physics:phyx-mini-0517:phyxminiproblems-problemphyxmini0517-actionquantity}
  \lean{PhyXMiniProblems.ProblemPhyXMini0517.ActionQuantity}
  \uses{def:physics:phyx-mini-0517:phyxminiproblems-problemphyxmini0517-actiondimension}
  A nonnegative, unit-independent physical value of action
\end{definition}
\begin{proof}
  This is the declared model definition of Action Quantity.
\end{proof}
\begin{definition}[Frequency Quantity]
  \label{def:physics:phyx-mini-0517:phyxminiproblems-problemphyxmini0517-frequencyquantity}
  \lean{PhyXMiniProblems.ProblemPhyXMini0517.FrequencyQuantity}
  A nonnegative, unit-independent ordinary frequency
\end{definition}
\begin{proof}
  This is the declared model definition of Frequency Quantity.
\end{proof}
\begin{definition}[energy In Joules]
  \label{def:physics:phyx-mini-0517:phyxminiproblems-problemphyxmini0517-energyinjoules}
  \lean{PhyXMiniProblems.ProblemPhyXMini0517.energyInJoules}
  Coherent-SI readout of an energy, in joules
\end{definition}
\begin{proof}
  This is the declared model definition of energy In Joules.
\end{proof}
\begin{definition}[energy In Electron Volts]
  \label{def:physics:phyx-mini-0517:phyxminiproblems-problemphyxmini0517-energyinelectronvolts}
  \lean{PhyXMiniProblems.ProblemPhyXMini0517.energyInElectronVolts}
  \uses{def:physics:phyx-mini-0517:phyxminiproblems-problemphyxmini0517-energyinjoules}
  Energy readout in electron volts, grounded by DimEnergy.electronVolt
\end{definition}
\begin{proof}
  This is the declared model definition of energy In Electron Volts.
\end{proof}
\begin{definition}[action In Joule Seconds]
  \label{def:physics:phyx-mini-0517:phyxminiproblems-problemphyxmini0517-actioninjouleseconds}
  \lean{PhyXMiniProblems.ProblemPhyXMini0517.actionInJouleSeconds}
  \uses{def:physics:phyx-mini-0517:phyxminiproblems-problemphyxmini0517-actionquantity}
  Coherent-SI readout of an action, in joule-seconds
\end{definition}
\begin{proof}
  This is the declared model definition of action In Joule Seconds.
\end{proof}
\begin{definition}[frequency In Hertz]
  \label{def:physics:phyx-mini-0517:phyxminiproblems-problemphyxmini0517-frequencyinhertz}
  \lean{PhyXMiniProblems.ProblemPhyXMini0517.frequencyInHertz}
  \uses{def:physics:phyx-mini-0517:phyxminiproblems-problemphyxmini0517-frequencyquantity}
  Coherent-SI readout of an ordinary frequency, in hertz
\end{definition}
\begin{proof}
  This is the declared model definition of frequency In Hertz.
\end{proof}
\begin{definition}[Atomic Energy Level]
  \label{def:physics:phyx-mini-0517:phyxminiproblems-problemphyxmini0517-atomicenergylevel}
  \lean{PhyXMiniProblems.ProblemPhyXMini0517.AtomicEnergyLevel}
  The three horizontal atomic energy levels in the diagram
\end{definition}
\begin{proof}
  This is the declared model definition of Atomic Energy Level.
\end{proof}
\begin{definition}[Emission Transition]
  \label{def:physics:phyx-mini-0517:phyxminiproblems-problemphyxmini0517-emissiontransition}
  \lean{PhyXMiniProblems.ProblemPhyXMini0517.EmissionTransition}
  The three downward arrows in the emission part of the diagram
\end{definition}
\begin{proof}
  This is the declared model definition of Emission Transition.
\end{proof}
\begin{definition}[initial Level]
  \label{def:physics:phyx-mini-0517:phyxminiproblems-problemphyxmini0517-emissiontransition-initiallevel}
  \lean{PhyXMiniProblems.ProblemPhyXMini0517.EmissionTransition.initialLevel}
  \uses{def:physics:phyx-mini-0517:phyxminiproblems-problemphyxmini0517-atomicenergylevel, def:physics:phyx-mini-0517:phyxminiproblems-problemphyxmini0517-emissiontransition}
  Initial level of each labeled emission transition
\end{definition}
\begin{proof}
  This is the declared model definition of initial Level.
\end{proof}
\begin{definition}[final Level]
  \label{def:physics:phyx-mini-0517:phyxminiproblems-problemphyxmini0517-emissiontransition-finallevel}
  \lean{PhyXMiniProblems.ProblemPhyXMini0517.EmissionTransition.finalLevel}
  \uses{def:physics:phyx-mini-0517:phyxminiproblems-problemphyxmini0517-atomicenergylevel, def:physics:phyx-mini-0517:phyxminiproblems-problemphyxmini0517-emissiontransition}
  Final level of each labeled emission transition
\end{definition}
\begin{proof}
  This is the declared model definition of final Level.
\end{proof}
\begin{definition}[Absorption Transition]
  \label{def:physics:phyx-mini-0517:phyxminiproblems-problemphyxmini0517-absorptiontransition}
  \lean{PhyXMiniProblems.ProblemPhyXMini0517.AbsorptionTransition}
  The two upward arrows in the absorption part of the diagram
\end{definition}
\begin{proof}
  This is the declared model definition of Absorption Transition.
\end{proof}
\begin{definition}[initial Level]
  \label{def:physics:phyx-mini-0517:phyxminiproblems-problemphyxmini0517-absorptiontransition-initiallevel}
  \lean{PhyXMiniProblems.ProblemPhyXMini0517.AbsorptionTransition.initialLevel}
  \uses{def:physics:phyx-mini-0517:phyxminiproblems-problemphyxmini0517-atomicenergylevel, def:physics:phyx-mini-0517:phyxminiproblems-problemphyxmini0517-absorptiontransition}
  Initial level of each labeled absorption transition
\end{definition}
\begin{proof}
  This is the declared model definition of initial Level.
\end{proof}
\begin{definition}[final Level]
  \label{def:physics:phyx-mini-0517:phyxminiproblems-problemphyxmini0517-absorptiontransition-finallevel}
  \lean{PhyXMiniProblems.ProblemPhyXMini0517.AbsorptionTransition.finalLevel}
  \uses{def:physics:phyx-mini-0517:phyxminiproblems-problemphyxmini0517-atomicenergylevel, def:physics:phyx-mini-0517:phyxminiproblems-problemphyxmini0517-absorptiontransition}
  Final level of each labeled absorption transition
\end{definition}
\begin{proof}
  This is the declared model definition of final Level.
\end{proof}
\begin{definition}[Arrow Direction]
  \label{def:physics:phyx-mini-0517:phyxminiproblems-problemphyxmini0517-arrowdirection}
  \lean{PhyXMiniProblems.ProblemPhyXMini0517.ArrowDirection}
  Vertical orientation of a transition arrow in the supplied diagram
\end{definition}
\begin{proof}
  This is the declared model definition of Arrow Direction.
\end{proof}
\begin{definition}[Arrow Color]
  \label{def:physics:phyx-mini-0517:phyxminiproblems-problemphyxmini0517-arrowcolor}
  \lean{PhyXMiniProblems.ProblemPhyXMini0517.ArrowColor}
  The two arrow colors which distinguish absorption from emission
\end{definition}
\begin{proof}
  This is the declared model definition of Arrow Color.
\end{proof}
\begin{definition}[Atomic Energy Level Diagram]
  \label{def:physics:phyx-mini-0517:phyxminiproblems-problemphyxmini0517-atomicenergyleveldiagram}
  \lean{PhyXMiniProblems.ProblemPhyXMini0517.AtomicEnergyLevelDiagram}
  \uses{def:physics:phyx-mini-0517:phyxminiproblems-problemphyxmini0517-atomicenergylevel, def:physics:phyx-mini-0517:phyxminiproblems-problemphyxmini0517-emissiontransition, def:physics:phyx-mini-0517:phyxminiproblems-problemphyxmini0517-absorptiontransition, def:physics:phyx-mini-0517:phyxminiproblems-problemphyxmini0517-arrowdirection, def:physics:phyx-mini-0517:phyxminiproblems-problemphyxmini0517-arrowcolor}
  Raster-visible diagram data. Its scalar labels are electron-volt readouts, not replacement types for the physical level energies stored in the setup
\end{definition}
\begin{proof}
  This is the declared model definition of Atomic Energy Level Diagram.
\end{proof}
\begin{definition}[Atom Preparation]
  \label{def:physics:phyx-mini-0517:phyxminiproblems-problemphyxmini0517-atompreparation}
  \lean{PhyXMiniProblems.ProblemPhyXMini0517.AtomPreparation}
  The preparation state specified by the words "when excited"
\end{definition}
\begin{proof}
  This is the declared model definition of Atom Preparation.
\end{proof}
\begin{definition}[Three Level Atom Setup]
  \label{def:physics:phyx-mini-0517:phyxminiproblems-problemphyxmini0517-threelevelatomsetup}
  \lean{PhyXMiniProblems.ProblemPhyXMini0517.ThreeLevelAtomSetup}
  \uses{def:physics:phyx-mini-0517:phyxminiproblems-problemphyxmini0517-actionquantity, def:physics:phyx-mini-0517:phyxminiproblems-problemphyxmini0517-frequencyquantity, def:physics:phyx-mini-0517:phyxminiproblems-problemphyxmini0517-atomicenergylevel, def:physics:phyx-mini-0517:phyxminiproblems-problemphyxmini0517-emissiontransition, def:physics:phyx-mini-0517:phyxminiproblems-problemphyxmini0517-atomicenergyleveldiagram, def:physics:phyx-mini-0517:phyxminiproblems-problemphyxmini0517-atompreparation}
  The physical atom and its emitted spectral lines. Each line frequency is a dimensionful observable. The governing-law premise below, rather than a definition here, relates it to the appropriate energy-level difference
\end{definition}
\begin{proof}
  This is the declared model definition of Three Level Atom Setup.
\end{proof}
\begin{definition}[Matches Excited Atom Scenario]
  \label{def:physics:phyx-mini-0517:phyxminiproblems-problemphyxmini0517-matchesexcitedatomscenario}
  \lean{PhyXMiniProblems.ProblemPhyXMini0517.MatchesExcitedAtomScenario}
  \uses{def:physics:phyx-mini-0517:phyxminiproblems-problemphyxmini0517-threelevelatomsetup}
  The atom is in the excited preparation contemplated by the question
\end{definition}
\begin{proof}
  This is the declared model definition of Matches Excited Atom Scenario.
\end{proof}
\begin{definition}[Matches Supplied Energy Level Diagram]
  \label{def:physics:phyx-mini-0517:phyxminiproblems-problemphyxmini0517-matchessuppliedenergyleveldiagram}
  \lean{PhyXMiniProblems.ProblemPhyXMini0517.MatchesSuppliedEnergyLevelDiagram}
  \uses{def:physics:phyx-mini-0517:phyxminiproblems-problemphyxmini0517-energyinelectronvolts, def:physics:phyx-mini-0517:phyxminiproblems-problemphyxmini0517-threelevelatomsetup}
  Numerical energy-level data and every visible feature of image 517. These are source/figure readouts only: no emitted frequency, maximum line, or answer choice occurs in this premise
\end{definition}
\begin{proof}
  This is the declared model definition of Matches Supplied Energy Level Diagram.
\end{proof}
\begin{definition}[Has Physical Spectral Parameters]
  \label{def:physics:phyx-mini-0517:phyxminiproblems-problemphyxmini0517-hasphysicalspectralparameters}
  \lean{PhyXMiniProblems.ProblemPhyXMini0517.HasPhysicalSpectralParameters}
  \uses{def:physics:phyx-mini-0517:phyxminiproblems-problemphyxmini0517-energyinjoules, def:physics:phyx-mini-0517:phyxminiproblems-problemphyxmini0517-actioninjouleseconds, def:physics:phyx-mini-0517:phyxminiproblems-problemphyxmini0517-frequencyinhertz, def:physics:phyx-mini-0517:phyxminiproblems-problemphyxmini0517-threelevelatomsetup}
  Positivity and strict level ordering for the physical branch of the model
\end{definition}
\begin{proof}
  This is the declared model definition of Has Physical Spectral Parameters.
\end{proof}
\begin{definition}[Uses Standard Planck Constant]
  \label{def:physics:phyx-mini-0517:phyxminiproblems-problemphyxmini0517-usesstandardplanckconstant}
  \lean{PhyXMiniProblems.ProblemPhyXMini0517.UsesStandardPlanckConstant}
  \uses{def:physics:phyx-mini-0517:phyxminiproblems-problemphyxmini0517-actioninjouleseconds, def:physics:phyx-mini-0517:phyxminiproblems-problemphyxmini0517-threelevelatomsetup}
  Physlib supplies the reduced Planck constant Constants.ℏ in SI joule-seconds. Ordinary Planck's constant is 2πℏ. This reference datum contains no line frequency or answer-choice value
\end{definition}
\begin{proof}
  This is the declared model definition of Uses Standard Planck Constant.
\end{proof}
\begin{definition}[Satisfies Planck Einstein Emission Law]
  \label{def:physics:phyx-mini-0517:phyxminiproblems-problemphyxmini0517-satisfiesplanckeinsteinemissionlaw}
  \lean{PhyXMiniProblems.ProblemPhyXMini0517.SatisfiesPlanckEinsteinEmissionLaw}
  \uses{def:physics:phyx-mini-0517:phyxminiproblems-problemphyxmini0517-energyinjoules, def:physics:phyx-mini-0517:phyxminiproblems-problemphyxmini0517-actioninjouleseconds, def:physics:phyx-mini-0517:phyxminiproblems-problemphyxmini0517-frequencyinhertz, def:physics:phyx-mini-0517:phyxminiproblems-problemphyxmini0517-emissiontransition, def:physics:phyx-mini-0517:phyxminiproblems-problemphyxmini0517-threelevelatomsetup}
  Planck--Einstein emission law E\_initial - E\_final = h ν for every arrow. It is a general relation between the independent physical quantities and does not specialize any frequency to the requested maximum or displayed answer
\end{definition}
\begin{proof}
  This is the declared model definition of Satisfies Planck Einstein Emission Law.
\end{proof}
\begin{definition}[emission Energy Gap In Joules]
  \label{def:physics:phyx-mini-0517:phyxminiproblems-problemphyxmini0517-emissionenergygapinjoules}
  \lean{PhyXMiniProblems.ProblemPhyXMini0517.emissionEnergyGapInJoules}
  \uses{def:physics:phyx-mini-0517:phyxminiproblems-problemphyxmini0517-energyinjoules, def:physics:phyx-mini-0517:phyxminiproblems-problemphyxmini0517-emissiontransition, def:physics:phyx-mini-0517:phyxminiproblems-problemphyxmini0517-threelevelatomsetup}
  Energy released by an emission transition, read in joules
\end{definition}
\begin{proof}
  This is the declared model definition of emission Energy Gap In Joules.
\end{proof}
\begin{lemma}[three To Ground has unique largest energy Gap]
  \label{lem:physics:phyx-mini-0517:phyxminiproblems-problemphyxmini0517-threetoground-has-unique-largest-energygap}
  \lean{PhyXMiniProblems.ProblemPhyXMini0517.threeToGround_has_unique_largest_energyGap}
  \uses{def:physics:phyx-mini-0517:phyxminiproblems-problemphyxmini0517-emissiontransition, def:physics:phyx-mini-0517:phyxminiproblems-problemphyxmini0517-threelevelatomsetup, def:physics:phyx-mini-0517:phyxminiproblems-problemphyxmini0517-matchessuppliedenergyleveldiagram, def:physics:phyx-mini-0517:phyxminiproblems-problemphyxmini0517-emissionenergygapinjoules}
  The direct 3 eV → 0 eV arrow has the unique largest energy gap
\end{lemma}
\begin{proof}
  The conclusion follows from the governing relations and figure readouts named in the statement of three To Ground has unique largest energy Gap.
\end{proof}
\begin{definition}[Is Unique Highest Emission Transition]
  \label{def:physics:phyx-mini-0517:phyxminiproblems-problemphyxmini0517-isuniquehighestemissiontransition}
  \lean{PhyXMiniProblems.ProblemPhyXMini0517.IsUniqueHighestEmissionTransition}
  \uses{def:physics:phyx-mini-0517:phyxminiproblems-problemphyxmini0517-frequencyinhertz, def:physics:phyx-mini-0517:phyxminiproblems-problemphyxmini0517-emissiontransition, def:physics:phyx-mini-0517:phyxminiproblems-problemphyxmini0517-threelevelatomsetup}
  A labeled line has strictly greater frequency than every other line
\end{definition}
\begin{proof}
  This is the declared model definition of Is Unique Highest Emission Transition.
\end{proof}
\begin{definition}[Answer Choice]
  \label{def:physics:phyx-mini-0517:phyxminiproblems-problemphyxmini0517-answerchoice}
  \lean{PhyXMiniProblems.ProblemPhyXMini0517.AnswerChoice}
  Labels of the four displayed multiple-choice answers
\end{definition}
\begin{proof}
  This is the declared model definition of Answer Choice.
\end{proof}
\begin{definition}[frequency Hertz]
  \label{def:physics:phyx-mini-0517:phyxminiproblems-problemphyxmini0517-answerchoice-frequencyhertz}
  \lean{PhyXMiniProblems.ProblemPhyXMini0517.AnswerChoice.frequencyHertz}
  \uses{def:physics:phyx-mini-0517:phyxminiproblems-problemphyxmini0517-answerchoice}
  Frequency printed beside an answer choice, in hertz
\end{definition}
\begin{proof}
  This is the declared model definition of frequency Hertz.
\end{proof}
\begin{definition}[recorded Dataset Answer]
  \label{def:physics:phyx-mini-0517:phyxminiproblems-problemphyxmini0517-recordeddatasetanswer}
  \lean{PhyXMiniProblems.ProblemPhyXMini0517.recordedDatasetAnswer}
  \uses{def:physics:phyx-mini-0517:phyxminiproblems-problemphyxmini0517-answerchoice}
  Dataset metadata recording answer D; this is not a theorem premise
\end{definition}
\begin{proof}
  This is the declared model definition of recorded Dataset Answer.
\end{proof}
\begin{definition}[answer Choice Error Hertz]
  \label{def:physics:phyx-mini-0517:phyxminiproblems-problemphyxmini0517-answerchoiceerrorhertz}
  \lean{PhyXMiniProblems.ProblemPhyXMini0517.answerChoiceErrorHertz}
  \uses{def:physics:phyx-mini-0517:phyxminiproblems-problemphyxmini0517-frequencyinhertz, def:physics:phyx-mini-0517:phyxminiproblems-problemphyxmini0517-emissiontransition, def:physics:phyx-mini-0517:phyxminiproblems-problemphyxmini0517-threelevelatomsetup, def:physics:phyx-mini-0517:phyxminiproblems-problemphyxmini0517-answerchoice}
  Absolute error between an emitted line and a displayed answer, in hertz
\end{definition}
\begin{proof}
  This is the declared model definition of answer Choice Error Hertz.
\end{proof}
\begin{definition}[Agrees With Displayed Frequency]
  \label{def:physics:phyx-mini-0517:phyxminiproblems-problemphyxmini0517-agreeswithdisplayedfrequency}
  \lean{PhyXMiniProblems.ProblemPhyXMini0517.AgreesWithDisplayedFrequency}
  \uses{def:physics:phyx-mini-0517:phyxminiproblems-problemphyxmini0517-emissiontransition, def:physics:phyx-mini-0517:phyxminiproblems-problemphyxmini0517-threelevelatomsetup, def:physics:phyx-mini-0517:phyxminiproblems-problemphyxmini0517-answerchoice, def:physics:phyx-mini-0517:phyxminiproblems-problemphyxmini0517-answerchoiceerrorhertz}
  Agreement after rounding the coefficient to the displayed hundredth in units of 10\textasciicircum{}14 Hz; half of one display unit is 5 × 10\textasciicircum{}11 Hz
\end{definition}
\begin{proof}
  This is the declared model definition of Agrees With Displayed Frequency.
\end{proof}
\begin{definition}[Is Unique Closest Answer Choice]
  \label{def:physics:phyx-mini-0517:phyxminiproblems-problemphyxmini0517-isuniqueclosestanswerchoice}
  \lean{PhyXMiniProblems.ProblemPhyXMini0517.IsUniqueClosestAnswerChoice}
  \uses{def:physics:phyx-mini-0517:phyxminiproblems-problemphyxmini0517-emissiontransition, def:physics:phyx-mini-0517:phyxminiproblems-problemphyxmini0517-threelevelatomsetup, def:physics:phyx-mini-0517:phyxminiproblems-problemphyxmini0517-answerchoice, def:physics:phyx-mini-0517:phyxminiproblems-problemphyxmini0517-answerchoiceerrorhertz}
  The named answer is strictly closer than every other displayed frequency
\end{definition}
\begin{proof}
  This is the declared model definition of Is Unique Closest Answer Choice.
\end{proof}
% --- Archon declaration topology end ---
% --- Archon physics formalization source end ---
